\subsection{Quasi-Newton Methods}

% Paleta de curvas de nivel
\colorlet{niv1}{AmarilloAviso!85!black}
\colorlet{niv2}{AmarilloAviso!50!VerdeNeon!75!black}
\colorlet{niv3}{VerdeNeon!65!black}
\colorlet{niv4}{VerdeNeon!40!CelesteBrillante!65!black}
\colorlet{niv5}{CelesteBrillante!65!black}
\colorlet{niv6}{CelesteBrillante!48!black}
\colorlet{niv7}{CelesteBrillante!34!black}
\providecommand{\reflibro}[1]{\par\vfill{\tiny\textcolor{GrisTexto!70}{\textit{Ref. libro:} #1}}}

% Comandos para las curvas de nivel de Quasi-Newton con gradiente de color
\newcommand{\qnValle}{%
    \foreach \r/\cl in {0.5/niv1, 1.0/niv2, 1.5/niv3, 2.0/niv4, 2.5/niv5, 3.0/niv6}
        \draw[\cl, thick, rotate around={35:(0,0)}] (0,0) ellipse ({\r*1.4} and {\r*0.4});%
}

\newcommand{\qnValleTrayectoria}{%
    \foreach \r/\cl in {0.3/niv1, 0.7/niv2, 1.2/niv3, 1.8/niv4, 2.5/niv5, 3.3/niv6}
        \draw[\cl, thick, rotate around={25:(0,0)}] (0,0) ellipse ({\r*1.6} and {\r*0.35});%
}

% --- FRAME 1: Introducción y Concepto ---
\begin{frame}{Quasi-Newton: Explicación Gráfica (1/3)}
    \begin{columns}[T]
        \begin{column}{0.44\textwidth}
            \centering
            \vspace{0.2cm}
            \begin{tikzpicture}[x=1.0cm, y=1.0cm, font=\scriptsize]
                \begin{scope}
                    \clip (-2.2,-2.2) rectangle (2.2,2.2);
                    
                    % Curvas de nivel con gradiente
                    \qnValle
                    
                    % Óptimo
                    \fill[white] (0,0) circle (1.5pt);
                    \draw[GrisTexto] (0,0) circle (1.5pt) node[below right, inner sep=3pt] {$x^*$};
                    
                    % Punto de inicio
                    \coordinate (X1) at (-1.5, 1.2);
                    \fill[GrisTexto!70] (X1) circle (1.6pt) node[above left, inner sep=2pt] {$x^{(k)}$};
                    
                    % Vector Gradiente Puro
                    \draw[CelesteBrillante, thick, dashed, -stealth] (X1) -- ++(0.7, 1.2) node[right, inner sep=2pt] {$-g^{(k)}$};
                    
                    % Vector Quasi-Newton
                    \draw[NaranjaBase, very thick, -stealth] (X1) -- (-0.4, 0.3) node[midway, below right, inner sep=1pt] {paso QN};
                    
                \end{scope}
                \draw[gray!50] (-2.2,-2.2) rectangle (2.2,2.2);
            \end{tikzpicture}
        \end{column}
        
        \begin{column}{0.54\textwidth}
            \scriptsize
            \begin{block}{De Newton a Quasi-Newton}
                El método de Newton exige el Hessiano exacto, calculando el paso como $x^{(k+1)} = x^{(k)} - (H^{(k)})^{-1} g^{(k)}$. En altas dimensiones, calcular e invertir esta matriz es prohibitivo.
            \end{block}
            
            \vspace{0.3cm}
            \begin{itemize}
                \setlength{\itemsep}{0.8em}
                \item \textbf{Estrategia QN:} Los métodos Quasi-Newton aproximan la inversa del Hessiano iterativamente.
                \item La regla de actualización se basa en una matriz $Q^{(k)}$ que aproxima dicha inversa en el punto $x^{(k)}$:
                \[ x^{(k+1)} \leftarrow x^{(k)} - \alpha^{(k)} Q^{(k)} g^{(k)} \]
            \end{itemize}
        \end{column}
    \end{columns}
\end{frame}

% --- FRAME 2: Descomposición del Paso (con la llave) ---
\begin{frame}{Quasi-Newton: Explicación Gráfica (2/3)}
    
    % --- FILA SUPERIOR: Ecuación y Gráfica ---
    \begin{columns}[T]
        \begin{column}{0.42\textwidth}
            \begin{block}{Iteración $k$}
                \footnotesize
                \[ \Delta x^{(k)} = -\alpha^{(k)} Q^{(k)} g^{(k)} \]
            \end{block}
        \end{column}
        
        \begin{column}{0.56\textwidth}
            \centering
            \begin{tikzpicture}[x=0.95cm, y=0.95cm, font=\scriptsize]
                \begin{scope}
                    \clip (-2.2,-2.2) rectangle (2.2,2.2);
                    
                    % Curvas de nivel con gradiente
                    \qnValle
                    
                    \coordinate (Xk) at (-1.5, 1.2);
                    
                    % Descomposición de vectores
                    \draw[CelesteBrillante, very thick, -stealth] (Xk) -- (-0.8, 2.4) node[midway, above left, inner sep=2pt] {$-g^{(k)}$};
                    \draw[VerdeNeon, very thick, dotted, -stealth] (-0.8, 2.4) -- (-0.4, 0.3) node[midway, right, inner sep=2pt] {corrección $Q^{(k)}$};
                    \draw[NaranjaBase, ultra thick, -stealth] (Xk) -- (-0.4, 0.3);
                    
                    \fill[GrisTexto!70] (Xk) circle (1.6pt) node[left, inner sep=3pt] {$x^{(k)}$};
                    \fill[NaranjaBase] (-0.4, 0.3) circle (1.8pt) node[right, inner sep=3pt, text=GrisTexto] {$x^{(k+1)}$};
                \end{scope}
                \draw[gray!50] (-2.2,-2.2) rectangle (2.2,2.2);
            \end{tikzpicture}
        \end{column}
    \end{columns}
    
    \vspace{-0.2cm}
    
    % --- FILA INFERIOR: Textos explicativos alineados ---
    \begin{columns}[T]
        \begin{column}{0.4\textwidth}
            \scriptsize
            \begin{itemize}
                \setlength{\itemsep}{0.5em}
                \item \tikz[remember picture, overlay] \coordinate (item1top); \textcolor{CelesteBrillante}{$-g^{(k)}$ (Gradiente)}: Apunta a la máxima pendiente local, provocando rebote.
                \item \textcolor{VerdeNeon}{$Q^{(k)}$ (Aprox. Hessiano)}: Multiplica al gradiente y al factor escalar $\alpha^{(k)}$ para aproximar el paso de Newton.\tikz[remember picture, overlay] \coordinate (item2bot);
            \end{itemize}
        \end{column}
        
        \begin{column}{0.6\textwidth}
            \scriptsize
            \vspace{0.6cm}
            \begin{itemize}
                \item \tikz[remember picture, overlay] \coordinate (item3); \textcolor{NaranjaBase}{Resultado}: El paso resultante se reorienta de forma eficiente hacia el óptimo.
            \end{itemize}
        \end{column}
    \end{columns}
    
    % --- Llave conectando las columnas inferiores ---
    \begin{tikzpicture}[remember picture, overlay]
        \coordinate (brace_x) at ([xshift=-26pt]item3);
        \draw[decorate, decoration={brace, amplitude=6pt}, thick, GrisTexto] 
            ([yshift=2ex]brace_x |- item1top) -- ([yshift=-0.5ex]brace_x |- item2bot);
    \end{tikzpicture}
    
\end{frame}

% --- FRAME 3: Trayectorias y Convergencia ---
\begin{frame}{Quasi-Newton: Explicación Gráfica (3/3)}
    \begin{columns}[T]
        \begin{column}{0.38\textwidth}
            \centering
            % Gráfico de Trayectorias
            \begin{tikzpicture}[x=0.9cm, y=0.9cm, font=\scriptsize]
                \begin{scope}
                    \clip (-2.5,-2.2) rectangle (2.5,1.5);
                    
                    % Curvas de nivel con gradiente para la trayectoria
                    \qnValleTrayectoria
                    
                    \fill[white] (0,0) circle (1.5pt);
                    \draw[GrisTexto] (0,0) circle (1.5pt) node[right, inner sep=3pt] {$x^*$};
                    
                    % Trayectoria GD (Zig-Zag)
                    \draw[CelesteBrillante, thick] (-2.2, 1.0) -- (-1.8, -0.6) -- (-1.0, 0.4) -- (-0.5, -0.3) -- (-0.1, 0.2) -- (0,0);
                    \fill[CelesteBrillante] (-2.2, 1.0) circle (1.2pt) node[above right, text=GrisTexto] {$x^{(1)}$};
                    \foreach \p in {(-1.8, -0.6), (-1.0, 0.4), (-0.5, -0.3), (-0.1, 0.2)} { \fill[CelesteBrillante] \p circle (1pt); }
                    
                    % Trayectoria Quasi-Newton
                    \draw[NaranjaBase, very thick] (-2.2, 1.0) -- (-0.8, 0.2) -- (-0.1, 0.05) -- (0,0);
                    \foreach \p in {(-0.8, 0.2), (-0.1, 0.05)} { \fill[NaranjaBase] \p circle (1.5pt); }
                \end{scope}
                \draw[gray!50] (-2.5,-2.2) rectangle (2.5,1.5);
                
                % Leyenda integrada
                \node[text=NaranjaBase, right] at (-2.4, -2.5) {\textbf{---} Quasi-Newton};
                \node[text=CelesteBrillante, right] at (0.2, -2.5) {\textbf{---} GD};
            \end{tikzpicture}
            
            \vspace{0.2cm}
            % Gráfico de Barras Corregido
            \begin{tikzpicture}[font=\scriptsize]
                \node[text=GrisTexto, font=\footnotesize] at (2.2, 1.8) {Distancia al mínimo ($||x_k - x^*||$)};
                \draw[gray] (0,0) -- (4.5,0);
                
                % Iteraciones 1 a 6
                \foreach \x/\qn/\gd in {1/1.5/1.5, 2/0.7/1.3, 3/0.15/1.1, 4/0.02/0.9, 5/0/0.75, 6/0/0.6} {
                    \fill[CelesteBrillante] ({\x*0.65+0.15},0) rectangle ++(0.2,\gd);
                    \fill[NaranjaBase] ({\x*0.65-0.1},0) rectangle ++(0.2,\qn);
                    \node[below] at ({\x*0.65+0.12},0) {\x};
                }
                \node[below] at (2.2, -0.3) {iteración $k$};
            \end{tikzpicture}
        \end{column}

        \scriptsize
        \begin{column}{0.60\textwidth}
            \begin{block}{Actualización BFGS}
                Se definen los cambios $\gamma^{(k+1)} \equiv g^{(k+1)} - g^{(k)}$ y $\delta^{(k+1)} \equiv x^{(k+1)} - x^{(k)}$. La aproximación se enriquece iterativamente:
                \vspace{-0.1cm}
                \[ Q \leftarrow Q - \left(\frac{\delta\gamma^\top Q + Q\gamma\delta^\top}{\delta^\top\gamma}\right) + \left(1 + \frac{\gamma^\top Q\gamma}{\delta^\top\gamma}\right)\frac{\delta\delta^\top}{\delta^\top\gamma} \]
            \end{block}
            
            \vspace{0.2cm}
            \begin{itemize}
                \setlength{\itemsep}{0.8em}
                \item \textbf{Costo de Memoria:} Aunque acelera la convergencia drásticamente respecto al GD, la matriz densa $Q$ demanda grandes cantidades de memoria espacial.
                \item \textbf{L-BFGS:} El método L-BFGS fue diseñado para aproximar a BFGS reduciendo este impacto. La matriz inversa completa se sustituye almacenando únicamente los últimos $m$ valores de $\delta$ y $\gamma$.
            \end{itemize}
        \end{column}
    \end{columns}
\end{frame}

% ---------------------------------------------------------------
% 2. Pseudocódigo
% ---------------------------------------------------------------
\begin{frame}{Quasi-Newton (BFGS): Pseudocódigo}
    \scriptsize
    \begin{algorithmic}[1]
        \Require $f$, $\nabla f$, $x^{(1)}$, tolerancia $\epsilon$, $k_{\max}$
        \State $Q \gets I$ \Comment{Matriz identidad: aprox. inicial de la inversa de la Hessiana}
        \For{$k = 1, \dots, k_{\max}$}
            \State $g \gets \nabla f(x)$
            \State \textbf{if} $\lVert g \rVert < \epsilon$ \textbf{then break}
            \State $d \gets -Q g$ \Comment{Dirección de descenso casi-Newton (6.22)}
            \State $x_{\text{new}} \gets \text{line\_search}(f, x, d)$ \Comment{Búsqueda de línea a lo largo de $d$}
            \State $g_{\text{new}} \gets \nabla f(x_{\text{new}})$
            \State $\delta \gets x_{\text{new}} - x$ \Comment{(6.24)}
            \State $\gamma \gets g_{\text{new}} - g$ \Comment{(6.23)}
            \State \textbf{if} $\delta^\top \gamma \approx 0$ \textbf{then break} \Comment{Prevenir división por cero}
            \State $Q \gets Q - \frac{\delta\gamma^\top Q + Q\gamma\delta^\top}{\delta^\top\gamma} + \left(1 + \frac{\gamma^\top Q\gamma}{\delta^\top\gamma}\right)\frac{\delta\delta^\top}{\delta^\top\gamma}$ \Comment{Actualización BFGS (6.26)}
            \State $x \gets x_{\text{new}}$
        \EndFor
        \State \Return $x$
    \end{algorithmic}

    \vspace{1em}
    \scriptsize
    \begin{itemize}
        \item Aproxima la inversa de la Hessiana ($H^{-1}$) usando solo información de gradientes pasados.
        \item Mantiene $Q$ simétrica y definida positiva en cada paso.
        \item Existen otras variantes como DFP (Ec. 6.25), pero BFGS es el estándar más robusto.
    \end{itemize}
    \reflibro{Algoritmo~6.6, Ec.~(6.22)--(6.26).}
\end{frame}

% ---------------------------------------------------------------
% 3. Ejemplo paso a paso
% ---------------------------------------------------------------
\begin{frame}{Quasi-Newton (BFGS): Ejemplo Paso a Paso (1/3)}
    \begin{exampleblock}{Planteamiento}
        \scriptsize
        $f(x)=\tfrac12x_1^2+5x_2^2$, \quad $\nabla f=(x_1,\ 10x_2)$, \quad
        $x^{(1)}=(-4,\,1)$, \quad $\alpha=0.15$
    \end{exampleblock}

    \begin{columns}[T]
        \begin{column}{0.5\textwidth}
            \begin{block}{Inicialización}
                \scriptsize
                $Q^{(1)}=I=\begin{pmatrix} 1 & 0 \\ 0 & 1 \end{pmatrix}$
            \end{block}
            \begin{block}{Iteración $k=1$}
                \scriptsize
                $g^{(1)}=(-4,\ 10)$\\[2pt]
                $d^{(1)}=-Q^{(1)}g^{(1)}=(4,\ -10)$\\[2pt]
                $x^{(2)}=x^{(1)} + 0.15\, d^{(1)} = (-3.40,\ -0.50)$\\[2pt]
                $\mathbf{f(x^{(2)})=7.030}$
            \end{block}
            {\scriptsize \textbf{Nota:} Al ser $Q^{(1)}=I$, el primer paso es idéntico a usar Descenso de Gradiente (GD).\par}
        \end{column}
        \begin{column}{0.48\textwidth}
            \centering
            \begin{tikzpicture}[x=1.05cm, y=1.05cm, font=\scriptsize]
                \begin{scope}
                    \clip (-4.4,-1.3) rectangle (0.6,1.35);
                    \momElipses
                    \draw[NaranjaBase, very thick, -stealth] (-4,1) -- (-3.40,-0.50);
                \end{scope}
                \draw[gray!50] (-4.4,-1.3) rectangle (0.6,1.35);
                \fill[GrisTexto] (-4,1) circle (1.8pt) node[right, inner sep=3pt] {$x^{(1)}$};
                \fill[GrisTexto] (-3.40,-0.50) circle (1.8pt) node[right, inner sep=3pt] {$x^{(2)}$};
                \fill[white] (0,0) circle (1.6pt) node[above, inner sep=2pt] {$x^*$};
            \end{tikzpicture}
        \end{column}
    \end{columns}
\end{frame}

\begin{frame}{Quasi-Newton (BFGS): Ejemplo Paso a Paso (2/3)}
    \begin{columns}[T]
        \begin{column}{0.56\textwidth}
            \begin{block}{Iteración $k=2$}
                \scriptsize
                $g^{(2)}=(-3.40,\ -5.00)$\\[2pt]
                $\delta^{(2)}=x^{(2)}-x^{(1)}=(0.60,\ -1.50)$\\[2pt]
                $\gamma^{(2)}=g^{(2)}-g^{(1)}=(0.60,\ -15.00)$\\[2pt]
                Actualizamos $Q^{(2)}$ usando la Ec. (6.26)\\[2pt]
                $d^{(2)}=-Q^{(2)}g^{(2)} \approx (3.90,\ 0.52)$\\[2pt]
                $x^{(3)}=x^{(2)}+0.15\, d^{(2)} = (-2.81,\ -0.42)$\\[2pt]
                $\mathbf{f(x^{(3)})=4.851}$
            \end{block}
            \small
            \begin{itemize}
                \item $Q$ se moldea capturando la curvatura.
                \item La dirección de descenso se desvía del gradiente negativo y apunta de forma más directa hacia el mínimo.
            \end{itemize}
        \end{column}
        \begin{column}{0.42\textwidth}
            \centering
            \begin{tikzpicture}[x=1.45cm, y=1.45cm, font=\scriptsize]
                \begin{scope}
                    \clip (-4.25,-1.75) rectangle (-1.9,1.15);
                    \momElipses
                    \draw[GrisTexto!60, thick, -stealth] (-4,1) -- (-3.40,-0.50);
                    \draw[NaranjaBase, very thick, -stealth] (-3.40,-0.50) -- (-2.81,-0.42);
                \end{scope}
                \draw[gray!50] (-4.25,-1.75) rectangle (-1.9,1.15);
                \fill[GrisTexto] (-4,1) circle (1.6pt) node[right, inner sep=3pt] {$x^{(1)}$};
                \fill[GrisTexto] (-3.40,-0.50) circle (1.6pt) node[left, inner sep=2pt] {$x^{(2)}$};
                \fill[VerdeNeon] (-2.81,-0.42) circle (1.8pt) node[right, inner sep=2pt, text=GrisTexto] {$x^{(3)}$};
            \end{tikzpicture}
        \end{column}
    \end{columns}
\end{frame}

\begin{frame}{Quasi-Newton (BFGS): Ejemplo Paso a Paso (3/3)}
    \begin{columns}[T]
        \begin{column}{0.56\textwidth}
            \begin{block}{Iteración $k=3$}
                \scriptsize
                $g^{(3)}=(-2.81,\ -4.22)$\\[3pt]
                Actualizamos $Q^{(3)}$ \dots\\[3pt]
                $x^{(4)}=(-2.39,\ -0.36)$\\[3pt]
                $\mathbf{f(x^{(4)})=3.505}$
            \end{block}
            \begin{exampleblock}{Resultado final ($k_{\max}=60$)}
                \scriptsize
                $x^{(61)} = (-0.0002,\ 0.0000)$\\[3pt]
                $f(x^{(61)}) = 0.0000$
            \end{exampleblock}
        \end{column}
        \begin{column}{0.42\textwidth}

            \centering
            \begin{tikzpicture}[x=0.95cm, y=0.95cm, font=\scriptsize]
                \begin{scope}
                    \clip (-4.4,-1.15) rectangle (0.5,1.3);
                    \momElipses
                    \draw[NaranjaBase, very thick] plot coordinates {(-4.00,1.00) (-3.40,-0.50) (-2.81,-0.42) (-2.39,-0.36) (-2.03,-0.30) (-1.73,-0.26) (-1.47,-0.22)};
                    \foreach \p in {(-3.40,-0.50), (-2.81,-0.42), (-2.39,-0.36), (-2.03,-0.30), (-1.73,-0.26), (-1.47,-0.22)}
                        \fill[NaranjaBase] \p circle (1.4pt);
                \end{scope}
                \draw[gray!50] (-4.4,-1.15) rectangle (0.5,1.3);
                \fill[white] (-4,1) circle (1.6pt) node[right, inner sep=3pt, text=GrisTexto] {$x^{(1)}$};
                \fill[white] (0,0) circle (1.4pt);
                \draw[NaranjaBase, very thick] (-4.3,-1.45) -- (-3.9,-1.45) node[right, text=GrisTexto] {BFGS};
            \end{tikzpicture}

            {\tiny Primeras 7 iteraciones de BFGS.\par}
        \end{column}
    \end{columns}

    \vspace{0.3em}
    \centering
    {\scriptsize
    \begin{tabular}{l c c c c c c c}
        \toprule
        $k$ & 0 & 1 & 2 & 3 & 4 & 5 & 6 \\
        \midrule
        $f(x^{(k+1)})$ BFGS & 13.00 & 7.030 & 4.851 & 3.505 & 2.532 & 1.830 & 1.322 \\
        \bottomrule
    \end{tabular}\par}
\end{frame}

% ---------------------------------------------------------------
% 4. Código
% ---------------------------------------------------------------
\begin{frame}[fragile]{Quasi-Newton (BFGS): Código}
\begin{lstlisting}[language=Python]
import numpy as np

def bfgs(f, grad, x, alpha=0.1, k_max=100, eps_tol=1e-6):
    m = len(x)
    Q = np.eye(m)  # Inicializacion de Q^{(1)} = I
    for k in range(1, k_max + 1):
        g = grad(x)
        if np.linalg.norm(g) < eps_tol:
            break         
        d = -Q @ g                        # (6.22) Direccion
        x_new = x + alpha * d             # Paso (reemplaza line_search)
        g_new = grad(x_new)
        delta = x_new - x                 # (6.24)
        gamma = g_new - g                 # (6.23)
        denom = np.dot(delta, gamma)
        if denom == 0:
            break    
        # Actualizacion BFGS (6.26)
        term1 = (np.outer(delta, gamma) @ Q + Q @ np.outer(gamma, delta)) / denom
        term2 = (1 + np.dot(gamma, Q @ gamma) / denom) * (np.outer(delta, delta) / denom)
        Q = Q - term1 + term2
        x = x_new
    return x, k

# Ejemplo: f(x) = 0.5*x1^2 + 5*x2^2   ->   grad f = (x1, 10*x2)
f = lambda x: 0.5*x[0]**2 + 5*x[1]**2
grad = lambda x: np.array([x[0], 10*x[1]])
x0 = np.array([-4.0, 1.0])
bfgs(f, grad, x0, alpha=0.1, k_max=100)
\end{lstlisting}
\end{frame}